\section{Problem Definition and System Description}

\subsection{Process architecture and operational analysis}

The Miller Pain Treatment Center (MPTC) operates as an outpatient chronic pain clinic within the Eastern Hospital Outpatient Center. The primary process is the delivery of pain treatment to outpatient visitors, and the main flow unit is the patient moving through a sequence of clinical and administrative activates \parencite{pensumbok}. The secondary flow unit is the patient's medical record, which precedes and enables each clinical interaction. The main process can be divided into three sub-processes. The first one being patient intake, then clinical assessment and finally patient discharge. The systems entry boundary is triggered when the signs in at the front desk, and the exit boundary is crossed when the patient completes the check-out process and leaves the facility. The center handles three types of patients: New (30\%) require extensive clinical interaction and arrive 30 minutes early, Return (40\%) follow the same full path with shorter activity times and Follow-up (30\%) bypass resident-driven steps entirely, routing directly to the PA. Process flow is shown in \autoref{tab:process}.

\vspace{8pt}
The AMC context adds significant structural complexity by introducing four resident-driven activities shown in figure 1 of the article \parencite[Figure 1]{MillerPainTreatmentCenter}. These steps reflect the AMC's dual mandate as a care provider and medical education center, resulting in longer, resource-intensive process where the Attending is involved multiple sequential and parallel activities. From the customer perspective, patients primarily require minimal waiting times and high-quality clinical outcomes. However, these patients must be balanced against AMCs perspective, being that these requirements must coexist with the clinics educational mandate, which extends patient cycle times to facilitate clinical training. 

\subsubsection{Design approach}

Dr. Weems approach to managing the MPTC is best characterized as incremental improvement, rather than a full redesign of the process \parencite{pensumbok}. Instead of restructuring the fundamental architecture of the AMC, he suggests targeted measures within the existing process. He suggests pre-processing to be able to reduce Resident Review and Teach times, also suggest schedule optimization to better match appointments intervals to service time variability, also an early PSC start to eliminate front-desk delays for the first patient and finally a late-arrival policy adapted from his private practice experience. Each of these ideas change a specific step or parameter but without altering the overall process or the roles of staff that is involved. This is consistent with the incremental improvement paradigm described in chapter 1, where changes are made at the margin to improve performance within an existing design rather than questioning the design itself \parencite{pensumbok}. 

\subsubsection{Design enablers}

Following Chapter 3 \parencite[sec. 3.4]{pensumbok}, four design enablers shape the MPTC. Firstly, people are the key constraint. The Attending is a critical bottleneck shared across three mutually exclusive activities. While residents add capacity, they introduce variability due to competing duties. Conversely, the PA provides a parallel channel that decouples Follow-up patients from the main flow. Secondly information flow is manual and paper-based, which creates strict sequential dependencies. Dr. Weems' pre-processing proposal directly targets this by shifting information review offline. Third, technology plays a limited role, the absence of electronic health records or automated tools restricts information transfer speed and prevents real-time schedule adjustments. Finally, facilities present a structural constraint. With one of the four exam rooms dedicated to the PA, the remaining three rooms create a bottleneck during peak loads, as patients require simultaneous availability of a room and clinical staff. 

\subsubsection{Capacity and bottleneck analysis}

Applying the capacity analysis framework from chapter 5, the Attending is the clear bottleneck in the AMC system. To assess this, consider the demand placed on the attending during a four-hour session. From the schedule template shown in table 6 in the case \parencite[Table 6]{MillerPainTreatmentCenter}, each session contains 5 New, 7 Return and 5 Follow-up appointments, resulting in 17 patients. For new patients, the attending is required for both Teach, mean 7.85 minutes, and Resident and Attending, mean 12.45 minutes, giving an expected attending time of approximately 20.3 minutes per new patient. For return patients the equivalent figures are 5.17 and 9.24 minutes, resulting in 14.4 minutes. Lastly for follow-up patients, the attending is only required in approximately 50\% of the cases for a brief consultation of around 3 minutes, giving an expected contribution of 1.5 minutes per follow-up patient. The total expected attending demand per session is therefore approximately:
\[(5 \cdot 20.3) + (7 \cdot 14.4) + (5 \cdot 1.5) = 209.8\ \text{minutes}\]
This will result in above the 240 minutes available in a four-hour session when registration, transition times and overruns are considered. This calculation confirms a mismatch in capacity load at the Attending level, this even before accounting for the high variability in activity times. 

The deterministic 15-imunte scheduling interval severely aggravates this capacity mismatch. While new patients require an expected process time of roughly 50 minutes, the highly variable Attending step (standard deviation = 12.21 min) mathematically guarantees that early overrun cascade, inflating waiting times for all subsequent patients. This illustrates the relationship between high utilization, variability and waiting time described in chapter 5 \parencite[p. 168-169]{pensumbok}. Additionally, resident availability acts as a secondary capacity constraint, where any absences increases individual workloads, directly delaying the critical Teach and Resident \& Attending steps that feed primary bottleneck. 

\subsection{Opportunities for improvement}

\subsubsection{Process-oriented analysis of weakness}

Applying the process-oriented perspective from chapter 2 \parencite{pensumbok}, three structural weaknesses can be identified. First, the AMC flow contains extensive sequential dependencies involving the Attending, creating a compounding chain of waiting. Resident Review occurs while the patient already occupies an exam room, forcing preparation time into the patients cycle time as pure non-value-adding delay. Second, the 15-minute scheduling interval is mismatched with the actual service time variability, we see that the Attending step alone has a standard deviation exceeding 12 minutes, meaning any early overrun propagates forward through the entire session. Third, the 10\% no-show rate result in irrecoverable idle capacity, since cancellations are discovered too late to fill the slot, reducing throughput without reducing staff workload. 

\subsubsection{Improvement opportunities using design principles}

The design principles outlined in section 3.6 \cite[p. 100-101]{pensumbok} provide a structured basis for evaluating Dr. Weems proposed interventions and identifying additional opportunities. These workflow design principles directly motivate the pre-processing proposal. By assigning each patient to a resident the evening before the visit, the resident will be able to complete the case review offline, eliminating the Resident Review step during the clinic session. Also, the Attending can conduct some of the case-specific teaching offline, reducing the Teach step during the session. As Dr. Weems estimates pre-processing will cut the average time for both activities in half, if this happens it would reduce the unit load on the Attending, which is the primary bottleneck identified in section 3.1a. 

\subsubsection{Schedule work based on its critical characteristics and eliminating buffers}

This principle supports the case of schedule optimization. Scheduling all appointment types at 15-minute intervals fails to reflect the high variability in New patients service times. A redesigned template allocating longer slots to New patients and shorter slots to Follow-up patients would better match actual service time distributions. This is well suited for simulation evaluation without disrupting real clinic operations. Similarly, the principle of eliminating non-value-adding buffers motivates the early Patient Service Coordinator (PSC) proposal. Having a PSC available at 7.30 a.m. ensures the first patient enters the clinical flow on time, which is critical since any initial delay propagates forward through the entire session. 


\subsubsection{Design the process for the dominant flow, not for the expectations}

The late-arrival policy successfully implemented in the private practice, where late arrivals was cut in half within 12 months, address patient-driven schedule disruption by protecting the dominant flow from expectations \parencite[p. 99]{pensumbok}. Transferring this policy to the AMC introduces uncertainty due to the different patient demographic and institutional culture, but its impact on schedule adherence can be evaluated by varying the arrival time distribution in the simulation. 


\subsubsection{Examine process interactions to avoid suboptimization}

A further opportunity not explicitly proposed by Dr. Weems is reducing teaching time variability across attending physicians. Since Teach occupies the Attending, and through the Attending all downstream patients, high variability in this step has an outsized impact on system performance, a classic case of suboptimization when viewed in isolation \parencite[p. 99-100]{pensumbok}. Standardizing the teaching interaction or moving part of it offline through pre-processing could reduce this variability and improve schedule predictability. 
